\subsection{UX guidelines}

\subsubsection*{Current conventions}

\begin{itemize}
  \item Spanish (Mexico) interface copy, two-digit tabular ticket numbers, and canonical ``turno''/``mostrador'' labels.
  \item Montserrat font files are bundled through \Code{@fontsource}; icons are code-native SVG paths.
  \item Controls generally provide labels or accessible names, focus-visible outlines, disabled states, status/live regions, and a skip link.
  \item Dialogs use the native \Code{dialog} element and require a second step for cancellation.
  \item Responsive CSS adapts cashier and public layouts; reduced-motion media queries suppress transitions/animations.
  \item Connection loss remains visible and mutation controls are disabled while cached state stays on screen.
\end{itemize}

\subsubsection*{Guidelines for changes}

Preserve two independent identifiers (visible number and UUID), do not introduce preparation language, keep destructive cancellation confirmed, and ensure every successful-looking mutation is backed by a successful server response. New configuration must state whether it is local, persistent, per-browser, or broadcast.

\subsubsection*{Required validation}

Viewing distance, minimum contrast/font size, screen orientation, browser/TV compatibility, keyboard-only behavior, assistive-technology behavior, touch targets, Spanish terminology, and reduced-motion acceptance all require formal testing. Existing CSS is evidence of intent, not proof of standards conformance.

\subsubsection*{Historical redesign notes migrated}

The removed \Code{docs/UI-REDESIGN.md} described the Troyanos blue/white visual identity, glass surfaces, responsive layouts, public-display pagination, and local assets. Those claims were retained here only where the current source still supports them. Its statements that public audio and a coffee welcome image were operational are superseded by current code evidence.
